% GENERATED FILE. Edit canonical lessons, then run npm run generate:books.
% canonical-source-hash: fnv1a64:dd8191cb38ac609a
% canonical-lessons: PT-C469-teimoso, PT-C469-teimosia, PT-C469-teimar, PT-C469-avarento, PT-C469-sovina

\chapter{Teimosos e sovinas}
\label{ch:pt-teimosos-e-sovinas}
\hlchaptermodality{469}

\begin{chapteropening}
\textbf{By the end of this chapter:} \emph{I can describe stubborn and stingy people.}

Complete the last lesson of chapter 469: I can describe stubborn and stingy people.
\end{chapteropening}

\section[teimoso]{teimoso — stubborn}

\label{lesson:PT-C469-teimoso}



\begin{quote}
Before the new one: say the Portuguese for selfish, then the Portuguese for selfishness.
\end{quote}

\subsection*{You'll want to know: teimoso}
\textbf{teimoso} — \textquotedblleft{}stubborn\textquotedblright{}.

\begin{quote}
O meu avô é teimoso e nunca muda de ideias.
\end{quote}

\begin{quote}
\emph{My grandfather is stubborn and never changes his mind.}
\end{quote}

Of things too: \emph{uma tosse teimosa} — a cough that will not go away.

\begin{cousinweb}
\emph{teimoso} is built on \emph{teima} \textquotedblleft{}obstinacy\textquotedblright{}, which is usually traced to Latin \emph{thema} \textquotedblleft{}a theme, a topic\textquotedblright{}, from Greek (EtymDB); the path of meaning is not fully settled.
\end{cousinweb}

\begin{grammarlens}[title={What you've built}]
One more word for telling, explaining and describing at B1.
\end{grammarlens}

\subsection*{Your turn}
\begin{itemize}
\raggedright
  \item \emph{Say it:} \emph{teimoso}
  \item \emph{Say it:} \emph{teimoso}, once more
  \item \emph{Say it:} a sentence of your own with \emph{teimoso}
  \item \emph{Recall:} say the Portuguese for selfish, then the Portuguese for selfishness, then say \emph{teimoso} again
\end{itemize}

\begin{tcolorbox}[breakable,colback=teal!4,colframe=teal!35!black,title={Before you move on}]
What does teimoso mean? (\textquotedblleft{}Stubborn\textquotedblright{}.) Say it once more.
\end{tcolorbox}
\section[a teimosia]{a teimosia — stubbornness}

\label{lesson:PT-C469-teimosia}



\begin{quote}
Before the new one: say the Portuguese for selfishness, then the Portuguese for stubborn.
\end{quote}

\subsection*{You'll want to know: a teimosia}
\textbf{a teimosia} — \textquotedblleft{}stubbornness\textquotedblright{}.

\begin{quote}
Por teimosia, ele não quis pedir ajuda.
\end{quote}

\begin{quote}
\emph{Out of stubbornness, he did not want to ask for help.}
\end{quote}

Built on \emph{teimoso}.

\begin{cousinweb}
\emph{teimosia} is built on \emph{teimoso} \textquotedblleft{}stubborn\textquotedblright{}, from \emph{teima} \textquotedblleft{}obstinacy\textquotedblright{}, which is usually traced to Latin \emph{thema} \textquotedblleft{}a theme\textquotedblright{} (EtymDB); that origin is not fully settled.
\end{cousinweb}

\begin{grammarlens}[title={What you've built}]
One more word for telling, explaining and describing at B1.
\end{grammarlens}

\subsection*{Your turn}
\begin{itemize}
\raggedright
  \item \emph{Say it:} \emph{a teimosia}
  \item \emph{Say it:} \emph{a teimosia}, once more
  \item \emph{Say it:} a sentence of your own with \emph{a teimosia}
  \item \emph{Recall:} say the Portuguese for selfishness, then the Portuguese for stubborn, then say \emph{a teimosia} again
\end{itemize}

\begin{tcolorbox}[breakable,colback=teal!4,colframe=teal!35!black,title={Before you move on}]
What does a teimosia mean? (\textquotedblleft{}Stubbornness\textquotedblright{}.) Say it once more.
\end{tcolorbox}
\section[teimar]{teimar — to insist (stubbornly)}

\label{lesson:PT-C469-teimar}



\begin{quote}
Before the new one: say the Portuguese for stubborn, then the Portuguese for stubbornness.
\end{quote}

\subsection*{You'll want to know: teimar}
\textbf{teimar} — \textquotedblleft{}to insist (stubbornly)\textquotedblright{}.

\begin{quote}
O miúdo teima em não comer a sopa.
\end{quote}

\begin{quote}
\emph{The kid insists on not eating his soup.}
\end{quote}

\emph{Teimar em} + infinitive — to insist on doing something.

\begin{cousinweb}
\emph{teimar} is built on \emph{teima} \textquotedblleft{}obstinacy\textquotedblright{}, usually traced to Latin \emph{thema} \textquotedblleft{}a theme\textquotedblright{}, from Greek (EtymDB); the origin is not fully settled.
\end{cousinweb}

\begin{grammarlens}[title={What you've built}]
One more word for telling, explaining and describing at B1.
\end{grammarlens}

\subsection*{Your turn}
\begin{itemize}
\raggedright
  \item \emph{Say it:} \emph{teimar}
  \item \emph{Say it:} \emph{teimar}, once more
  \item \emph{Say it:} a sentence of your own with \emph{teimar}
  \item \emph{Recall:} say the Portuguese for stubborn, then the Portuguese for stubbornness, then say \emph{teimar} again
\end{itemize}

\begin{tcolorbox}[breakable,colback=teal!4,colframe=teal!35!black,title={Before you move on}]
What does teimar mean? (\textquotedblleft{}To insist (stubbornly)\textquotedblright{}.) Say it once more.
\end{tcolorbox}
\section[avarento]{avarento — miserly; miser}

\label{lesson:PT-C469-avarento}



\begin{quote}
Before the new one: say the Portuguese for stubbornness, then the Portuguese for to insist.
\end{quote}

\subsection*{You'll want to know: avarento}
\textbf{avarento} — \textquotedblleft{}miserly; miser\textquotedblright{}.

\begin{quote}
O velho avarento guardava o dinheiro debaixo da cama.
\end{quote}

\begin{quote}
\emph{The miserly old man kept his money under the bed.}
\end{quote}

The vice is \emph{a avareza}.

\begin{cousinweb}
\emph{avarento} is built on \emph{avaro} \textquotedblleft{}greedy\textquotedblright{}, which goes back to Latin \emph{avārus} \textquotedblleft{}greedy, covetous\textquotedblright{} (EtymDB; Lewis \& Short).
\end{cousinweb}

\begin{grammarlens}[title={What you've built}]
One more word for telling, explaining and describing at B1.
\end{grammarlens}

\subsection*{Your turn}
\begin{itemize}
\raggedright
  \item \emph{Say it:} \emph{avarento}
  \item \emph{Say it:} \emph{avarento}, once more
  \item \emph{Say it:} a sentence of your own with \emph{avarento}
  \item \emph{Recall:} say the Portuguese for stubbornness, then the Portuguese for to insist, then say \emph{avarento} again
\end{itemize}

\begin{tcolorbox}[breakable,colback=teal!4,colframe=teal!35!black,title={Before you move on}]
What does avarento mean? (\textquotedblleft{}Miserly; miser\textquotedblright{}.) Say it once more.
\end{tcolorbox}
\section[sovina]{sovina — stingy; skinflint}

\label{lesson:PT-C469-sovina}



\begin{quote}
Before the new one: say the Portuguese for to insist, then the Portuguese for miserly.
\end{quote}

\subsection*{You'll want to know: sovina}
\textbf{sovina} — \textquotedblleft{}stingy; skinflint\textquotedblright{}.

\begin{quote}
O meu tio é sovina: nunca paga um café a ninguém.
\end{quote}

\begin{quote}
\emph{My uncle is stingy: he never buys anyone a coffee.}
\end{quote}

Informal. The same form for a man and a woman.

\begin{cousinweb}
The origin of \emph{sovina} is not settled; dictionaries offer several guesses and none is certain.
\end{cousinweb}

\begin{grammarlens}[title={What you've built}]
One more word for telling, explaining and describing at B1.
\end{grammarlens}

\subsection*{Your turn}
\begin{itemize}
\raggedright
  \item \emph{Say it:} \emph{sovina}
  \item \emph{Say it:} \emph{sovina}, once more
  \item \emph{Say it:} a sentence of your own with \emph{sovina}
  \item \emph{Recall:} say the Portuguese for to insist, then the Portuguese for miserly, then say \emph{sovina} again
\end{itemize}

\begin{tcolorbox}[breakable,colback=teal!4,colframe=teal!35!black,title={Before you move on}]
What does sovina mean? (\textquotedblleft{}Stingy; skinflint\textquotedblright{}.) Say it once more.
\end{tcolorbox}
